% What's Real and What Isn't: A Transparent Assessment of
% Consciousness-Physics Coupling in Agent Simulation
%
% Copyright (c) 2026 Tristan Stoltz / Luminous Dynamics
% SPDX-License-Identifier: AGPL-3.0-or-later
%
% Target: Minds and Machines

\documentclass[12pt,a4paper]{article}

\usepackage[utf8]{inputenc}
\usepackage[T1]{fontenc}
\usepackage{lmodern}
\usepackage[margin=1in]{geometry}
\usepackage{booktabs}
\usepackage{array}
\usepackage{tabularx}
\usepackage{graphicx}
\usepackage{amsmath}
\usepackage{amssymb}
\usepackage{natbib}
\usepackage{hyperref}
\usepackage{xcolor}
% \usepackage{enumitem} % not available on this system

\hypersetup{
  colorlinks=true,
  linkcolor=blue!60!black,
  citecolor=blue!60!black,
  urlcolor=blue!60!black
}

\newcolumntype{L}[1]{>{\raggedright\arraybackslash}p{#1}}

\title{What's Real and What Isn't:\\
A Transparent Assessment of Consciousness-Physics Coupling\\
in Agent Simulation}

\author{Tristan Stoltz\\
  Luminous Dynamics\\
  \texttt{tristan.stoltz@evolvingresonantcocreationism.com}}

\date{April 2026}

\begin{document}

\maketitle

% ============================================================
\begin{abstract}
A growing number of agent simulation engines couple internal
integration metrics---often labeled ``consciousness''---to
physical dynamics, claiming that an agent's integration state
influences its motor behavior, sensory processing, or even the
geometry of its simulated environment.  Most such systems do not
disclose which of their mechanics rest on established physics,
which draw loose inspiration from neuroscience, and which are
pure invention.  We propose a three-tier classification
framework---\textsc{Grounded}, \textsc{Inspired}, and
\textsc{Speculative}---and apply it, with full candor, to our
own engine (Symtropy).  Of nine distinct mechanics, we classify
two as \textsc{Grounded}, three as \textsc{Inspired}, and four
as \textsc{Speculative}.  We report what changed in our design
process once we performed this audit: speculative mechanics were
not removed but were relabeled, reframed, and documented as
computational metaphors rather than physical claims.  We argue
that this kind of transparency should be standard practice in
any system that couples integration metrics to simulated physics,
and we propose concrete reporting standards for the field.
\end{abstract}

\section{Introduction}
\label{sec:introduction}

The last decade has seen a proliferation of systems that attempt
to simulate, measure, or operationalize aspects of integrated
information processing in artificial agents.  Implementations of
Integrated Information Theory
\citep{tononi2004,oizumi2014} now appear in game AI, digital
twins, and safety-testing frameworks for autonomous systems.
Free Energy Principle (FEP) formulations
\citep{friston2010} are used to drive robot behavior.
Claims of ``conscious AI'' appear in both academic preprints and
press releases with increasing frequency.

Many of these systems couple an internal integration metric---a
scalar or vector quantity that summarizes how tightly an agent's
information processing is integrated---to the agent's physical
dynamics.  The coupling takes various forms: the metric may gate
motor output, modulate sensory gain, influence decision
thresholds, or, in more adventurous designs, bend the geometry
of the simulated environment itself.

The problem is not that these couplings exist.  Simulation
engines are free to define whatever mechanics serve their design
goals.  The problem is that most systems do not clearly
distinguish between three very different kinds of mechanics:

\begin{enumerate}
  \item[(i)] Mechanics that implement well-established physical or
    mathematical laws and can be validated against nature;
  \item[(ii)] Mechanics that draw loose analogy from real systems but
    do not faithfully implement their equations; and
  \item[(iii)] Mechanics that have no published precedent and are
    entirely invented for the simulation.
\end{enumerate}

When these categories are conflated, the result is what we call
\emph{integration washing}: the presentation of speculative
computational metaphors as though they were grounded in
established science.  This is not always intentional.
Developers who build engines iteratively may lose track of which
coupling channels rest on real physics and which rest on
aesthetic intuition.  But the downstream consequences are
serious.

For \textbf{AI safety}, reviewers evaluating an agent system
that claims ``integration-coupled physics'' need to know whether
the coupling is a validated thermodynamic constraint or an
invented force field.  The safety implications differ
categorically.

For \textbf{reproducibility}, speculative mechanics cannot be
validated against nature because there is no natural system they
model.  They can only be validated against their own
specification---a tautology that should be disclosed.

For \textbf{philosophy of mind}, conflating grounded and
speculative mechanics muddles the already difficult question of
whether any computational system implements genuine integration
\citep{chalmers1995}.  If a system claims to couple
``information integration to spacetime curvature'' but the
coupling is an invented scalar multiplier, the philosophical
import is very different from what the language implies.

For \textbf{education}, students encountering these systems
deserve to know which parts teach real physics and which parts
teach game design.

In this paper, we propose a classification framework, apply it
with full honesty to our own engine, report what the audit
changed, and argue that this transparency should be standard.


% ============================================================
\section{The Classification Framework}
\label{sec:framework}

We propose that every mechanic in an integration-coupled
simulation engine be classified into exactly one of three tiers.
The tiers are distinguished by the strength of the connection
between the mechanic and published, empirically validated
science.

\subsection{\textsc{Grounded}}

A mechanic is \textsc{Grounded} if it satisfies all of the
following:

\begin{enumerate}
  \item It implements equations drawn from a published physical
    or mathematical theory (with citation).
  \item The constants used are drawn from the literature, not
    invented.
  \item The mechanic's predictions, within its domain of
    application, can in principle be validated against physical
    measurement.
\end{enumerate}

Example: a thermodynamic subsystem that computes internal energy
$U$, temperature $T$, entropy $S$, and Helmholtz free energy
$F = U - TS$ using Landauer's principle \citep{landauer1961} for
the energy cost of bit erasure.  The equations are textbook
thermodynamics.  The constants ($k_B T \ln 2$ per erased bit)
are empirically established.  The predictions (minimum energy
dissipation per computation) are experimentally verified
\citep{berut2012}.

A \textsc{Grounded} classification does not mean the mechanic
is correctly applied to the simulation domain---only that the
underlying physics is real.  Whether thermodynamic constraints
meaningfully apply to an abstract agent's ``cognitive''
processes is a separate question.

\subsection{\textsc{Inspired}}

A mechanic is \textsc{Inspired} if:

\begin{enumerate}
  \item It draws on a real scientific theory or empirical
    finding, and the connection is stated explicitly.
  \item It captures the \emph{qualitative spirit} of the source
    theory but does not faithfully implement its equations,
    constants, or boundary conditions.
  \item The developer acknowledges the gap between the source
    theory and the implementation.
\end{enumerate}

Example: a free-energy gradient that drives agent behavior ``in
the spirit of'' the Free Energy Principle \citep{friston2010}
but uses a simplified scalar proxy rather than the full
variational inference framework.  The inspiration is legitimate;
the implementation is a loose analogy.

The key disclosure for an \textsc{Inspired} mechanic is:
\emph{what, specifically, was simplified, omitted, or
reinterpreted?}  Without this, the reader cannot assess how far
the implementation has drifted from its source.

\subsection{\textsc{Speculative}}

A mechanic is \textsc{Speculative} if:

\begin{enumerate}
  \item No published theory or empirical finding supports the
    specific coupling it implements.
  \item The mechanic is the developer's invention, even if it
    uses vocabulary borrowed from real physics.
  \item There is no known natural system in which this coupling
    has been observed.
\end{enumerate}

Example: a mechanic that couples an agent's integration metric
$\Phi$ to the curvature of its simulated spacetime, producing a
``$\Phi$-gravity'' effect.  Tononi's IIT
\citep{tononi2004} never proposed gravitational effects of
integrated information.  General relativity \citep{carroll2004}
couples the stress-energy tensor to curvature, not information
integration.  The mechanic borrows vocabulary from both theories
but implements a coupling that neither theory contains.

A \textsc{Speculative} classification is not a condemnation.
Speculative mechanics can produce interesting emergent behavior,
can serve as generative hypotheses, and can be valuable game
design tools.  The classification simply requires that they be
presented honestly: as inventions, not as implementations of
established science.


% ============================================================
\section{The Engine Under Audit}
\label{sec:engine}

Symtropy is an agent simulation engine with five coupling
channels between an internal integration metric and simulated
physics.  The engine implements nine distinct mechanics across
these channels.  It has been exercised across more than 3,000
simulation runs during development and testing.

The coupling channels are:

\begin{enumerate}
  \item \textbf{Thermodynamic}: The agent maintains an internal
    thermodynamic state ($U$, $T$, $S$, $F$) that constrains
    its computational budget.
  \item \textbf{Geometric}: The simulated environment uses
    conformal geometry, and certain mechanics modulate curvature
    or metric properties.
  \item \textbf{Free-energy}: A scalar gradient drives the agent
    toward states that minimize a prediction-error proxy.
  \item \textbf{Field}: Scalar fields in the environment are
    influenced by agent state or by spatial configuration.
  \item \textbf{Motor}: The agent's motor output is gated by
    its integration state.
\end{enumerate}

The engine was built iteratively over approximately eighteen
months.  During that period, mechanics were added based on a
mixture of scientific reading, design intuition, and aesthetic
preference.  The audit described in this paper was conducted
after the engine was functionally complete, precisely because the
developers (the present author included) realized they could no
longer clearly articulate which mechanics rested on real physics
and which did not.

This is, we suspect, a common situation.  It is also the
situation that our proposed framework is designed to address.


% ============================================================
\section{Application: The Realism Audit}
\label{sec:audit}

We applied the three-tier framework to each of Symtropy's nine
mechanics.  The results are summarized in
Table~\ref{tab:audit} and discussed individually below.

\begin{table}[!htbp]
\centering
\small
\caption{Realism audit of Symtropy's integration-physics
  coupling mechanics.  Each mechanic is classified as
  \textsc{Grounded} (G), \textsc{Inspired} (I), or
  \textsc{Speculative} (S).}
\label{tab:audit}
\smallskip
\begin{tabularx}{\textwidth}{L{3.2cm} c L{8.5cm}}
\toprule
\textbf{Mechanic} & \textbf{Tier} & \textbf{Justification} \\
\midrule
Thermodynamics ($U$, $T$, $S$, $F$)
  & G
  & Implements Landauer's principle \citep{landauer1961} and
    the second law.  Constants from literature.  Experimentally
    validated \citep{berut2012}. \\[6pt]
Conformal geometry (math)
  & G
  & Uses standard differential geometry
    \citep{carroll2004,wald1984}.  Equations are textbook.
    No novel physical claims. \\[6pt]
FEP gradient
  & I
  & Captures the qualitative drive toward prediction-error
    minimization \citep{friston2010} but uses a simplified
    scalar proxy, not the full variational framework.
    Omits generative model, recognition density, and
    Markov blanket formalism. \\[6pt]
Harmony fields
  & I
  & Loosely inspired by McFadden's CEMI field theory
    \citep{mcfadden2020}, which proposes that
    electromagnetic fields in the brain contribute to
    integration.  Our fields are scalar, not
    electromagnetic, and lack Maxwell's equations.  The
    qualitative idea (fields as integration medium) is
    preserved; the physics is not. \\[6pt]
Motor gain tiers
  & I
  & Inspired by motor cortex gating, where preparatory
    activity modulates the gain of motor output
    \citep{kaufman2014}.  Our implementation uses discrete
    tiers rather than continuous population dynamics.
    The qualitative structure (integration state gates
    motor authority) has biological precedent; the
    specific tier boundaries do not. \\[6pt]
Conformal curvature as physics
  & S
  & No published physical law couples an information
    integration metric to spacetime curvature.  The
    mathematics is valid (conformal rescaling of a
    metric), but the \emph{physical claim} that
    integration produces curvature is our invention.
    General relativity couples the stress-energy tensor
    to curvature \citep{carroll2004}, not information
    integration. \\[6pt]
Sanctuary zones
  & S
  & No precedent for spatial regions where a ``harmony
    field'' suppresses entropy or confers protective
    effects.  The mechanic is a game-design construct
    presented in physics vocabulary. \\[6pt]
$\Phi$-gravity
  & S
  & Tononi's IIT \citep{tononi2004} defines $\Phi$ as
    a measure of integrated information.  It does not
    propose gravitational effects.  Our coupling of
    $\Phi$ to a gravitational potential is pure
    invention. \\[6pt]
Dimensional leakage
  & S
  & Borrows the language of Randall-Sundrum brane
    scenarios \citep{randall1999}, but the actual
    mechanic (integration metric modulating
    cross-dimensional information flow) has no
    precedent in braneworld physics or any other
    published theory. \\
\bottomrule
\end{tabularx}
\end{table}

\subsection{Grounded Mechanics}

\paragraph{Thermodynamics.}
The engine's thermodynamic subsystem computes internal energy,
temperature, entropy, and Helmholtz free energy using standard
equations.  The minimum energy cost of computation is set by
Landauer's bound: $E_{\min} = k_B T \ln 2$ per bit erased
\citep{landauer1961}.  This has been experimentally verified to
within a factor of order unity \citep{berut2012}.

We note, however, that applying Landauer's bound to an abstract
agent's ``cognitive'' operations is itself a modeling choice.
The bound applies to physical bit erasure in thermodynamic
systems.  Whether it constrains computation in an abstract
simulation is a question about the simulation's semantics, not
about thermodynamics.  We classify the mechanic as
\textsc{Grounded} because the equations and constants are real,
while acknowledging that the \emph{domain of application} is
a design decision.

\paragraph{Conformal geometry.}
The engine uses conformal transformations---rescalings of a
metric tensor by a smooth positive function---as a mathematical
framework for environment geometry.  The mathematics is standard
\citep{carroll2004,wald1984}.  No novel physical claims are
made by the mathematical framework itself.  The \emph{use} of
this framework (what drives the conformal factor) is where
classifications diverge, as discussed under ``Conformal
curvature as physics'' below.

\subsection{Inspired Mechanics}

\paragraph{FEP gradient.}
The Free Energy Principle \citep{friston2010} proposes that
biological agents minimize variational free energy, which
requires a generative model, a recognition density, and a
Markov blanket.  Our implementation preserves the qualitative
idea---agents are driven toward states that reduce prediction
error---but replaces the full variational apparatus with a
scalar gradient computed from a simplified surprise signal.

What is lost: the rich structure of approximate Bayesian
inference, the distinction between perceptual and active
inference, and the information-geometric properties of the
variational bound.  What is preserved: the directional pressure
toward prediction-error reduction as a behavioral driver.

We classify this as \textsc{Inspired} because the simplification
is severe enough that the mechanic cannot be called an
implementation of the FEP.  It is a scalar proxy that shares the
FEP's motivating intuition.

\paragraph{Harmony fields.}
McFadden's Conscious Electromagnetic Information (CEMI) theory
\citep{mcfadden2020} proposes that the brain's endogenous
electromagnetic field is not an epiphenomenon but a functional
medium for information integration.  This is a minority position
in neuroscience, though it has generated testable predictions
\citep{mcfadden2020}.

Our ``harmony fields'' borrow the qualitative idea that a field
can serve as an integration medium.  However, our fields are
scalar (not electromagnetic), do not obey Maxwell's equations,
and are not coupled to any physical field dynamics.  The
name ``harmony field'' itself signals that we have moved from
physics to metaphor.

We considered classifying this as \textsc{Speculative}, since
the connection to CEMI is so loose.  We settled on
\textsc{Inspired} because the core idea (fields as integration
substrates) does have a published scientific advocate, even if
our implementation shares almost none of the theory's formal
content.

\paragraph{Motor gain tiers.}
Neurophysiology has established that motor cortex activity
undergoes preparatory modulation before movement execution:
population-level dynamics in premotor and primary motor cortex
gate the gain of downstream motor commands
\citep{kaufman2014}.  Our implementation captures this
qualitative structure---the agent's integration state determines
how much motor authority it has---but uses discrete tiers
(four levels) rather than continuous population dynamics.

The tier boundaries are not derived from neurophysiology.  They
are design parameters tuned for gameplay.  The classification is
\textsc{Inspired}: the structure has biological precedent, but
the parameterization does not.

\subsection{Speculative Mechanics}

\paragraph{Conformal curvature as physics.}
The engine allows the agent's integration metric to drive the
conformal factor of the environment's geometry.  This means that
an agent's integration state literally bends the simulated
space around it.

This is our most problematic mechanic from a transparency
standpoint.  The mathematical framework (conformal geometry) is
\textsc{Grounded}.  The \emph{physical claim} that integration
produces curvature is \textsc{Speculative}.  There is no
published physical law that couples any information-theoretic
quantity to spacetime curvature.  General relativity couples the
stress-energy tensor to the Einstein tensor
\citep{carroll2004}.  IIT defines $\Phi$ as a property of
causal structures \citep{tononi2004}.  No bridge between these
two theories has been established.

The mechanic produces visually compelling emergent behavior:
high-integration agents warp the local geometry, creating
navigational consequences for nearby agents.  This is
interesting \emph{as game design}.  It is not physics.

\paragraph{Sanctuary zones.}
The engine designates certain spatial regions as ``sanctuaries''
where harmony field values are elevated and entropy production
is suppressed.  There is no physical or biological precedent for
spatial regions that confer integrative protection by virtue of
their location.

The mechanic serves a game-design function: it creates
strategically important locations that reward integration-aware
play.  We classify it as \textsc{Speculative} because it is
presented using physics vocabulary (entropy suppression, field
elevation) without any physical basis.

\paragraph{$\Phi$-gravity.}
The engine implements a gravitational-like attraction between
agents that is proportional to their integration metrics.  This
directly borrows the language and mathematical form of
Newtonian gravitation but replaces mass with $\Phi$.

Tononi's IIT \citep{tononi2004,tononi2016} defines $\Phi$ as
the irreducible integrated information of a system.  The theory
makes no claims about gravitational effects.  Our coupling is
pure invention.  It produces interesting flocking and clustering
behavior, but it is not a physical prediction of IIT or any
other published theory.

\paragraph{Dimensional leakage.}
The engine includes a mechanic where high integration values
allow information to ``leak'' between simulated dimensions,
borrowing the language of Randall-Sundrum braneworld scenarios
\citep{randall1999}.  In the original Randall-Sundrum models,
gravitational fields propagate in a five-dimensional bulk while
Standard Model fields are confined to a four-dimensional brane.
The mechanism for confinement involves exponential warp factors
in the bulk metric.

Our mechanic preserves none of this structure.  The
``dimensions'' in our engine are abstract game-state dimensions,
not spatial dimensions.  The ``leakage'' is controlled by the
integration metric, not by bulk geometry.  The Randall-Sundrum
vocabulary is borrowed for evocative effect, not for physical
content.

We flag this as the clearest case of vocabulary borrowing that
could mislead: a reader encountering ``dimensional leakage
inspired by Randall-Sundrum'' might reasonably assume a
connection to extra-dimensional physics that does not exist.


% ============================================================
\section{What the Audit Changed}
\label{sec:lessons}

Performing this audit produced five concrete changes in our
development practice.

\paragraph{1. Terminology.}
Before the audit, internal documentation referred to ``consciousness
bends space'' as a description of the conformal curvature
mechanic.  After the audit, we reframed this as
``integration-parameterized curvature.''  The mechanic is
identical; the description is honest.  We found that the revised
language also improved technical communication: ``integration
metric'' has a clear referent (a scalar computed from agent
state), while ``consciousness'' imports philosophical baggage
that obscures the engineering.

\paragraph{2. Documentation.}
Each mechanic now carries a classification tag in its source
code documentation.  Developers encountering a mechanic for
the first time can immediately see whether it rests on real
physics, loose analogy, or invention.

\paragraph{3. Interpretation of results.}
Emergent behaviors arising from \textsc{Speculative} mechanics
are now reported differently.  Before the audit, we might have
written: ``the engine demonstrates that integration produces
gravitational effects.''  After the audit, we write: ``the
engine's invented $\Phi$-gravity mechanic produces clustering
behavior that is interesting as a simulation artifact.''  The
second formulation is less exciting but more honest.

\paragraph{4. No mechanics were removed.}
This is important.  The audit did not lead us to strip out
speculative mechanics.  They serve design goals, produce
interesting behavior, and may even generate testable hypotheses.
What changed is how they are \emph{presented}.  A
\textsc{Speculative} mechanic clearly labeled as such is a tool.
A \textsc{Speculative} mechanic presented as \textsc{Grounded}
is a misrepresentation.

\paragraph{5. Humility about ``Inspired'' mechanics.}
The audit revealed that our \textsc{Inspired} mechanics were
the hardest to classify.  The FEP gradient ``captures the
spirit'' of the Free Energy Principle---but how much spirit is
enough to justify the label?  We adopted a conservative rule:
if the implementation omits the core mathematical structure of
the source theory (not just incidental details), the mechanic
should be classified as \textsc{Inspired} at best, regardless
of how faithfully it captures the qualitative intuition.


% ============================================================
\section{Proposed Standards}
\label{sec:standards}

Based on our experience, we propose that the following
disclosures become standard practice for any system that couples
an integration metric (or any quantity labeled ``consciousness,''
``awareness,'' ``sentience,'' etc.) to simulated physics.

\subsection{The Realism Table}

Every paper describing such a system should include a table,
analogous to Table~\ref{tab:audit}, that classifies each
coupling mechanic as \textsc{Grounded}, \textsc{Inspired}, or
\textsc{Speculative}.  The table should include:

\begin{itemize}
  \item The name of the mechanic.
  \item The classification tier.
  \item A brief justification citing the relevant published
    science (for \textsc{Grounded} and \textsc{Inspired}) or
    explicitly stating the absence of precedent (for
    \textsc{Speculative}).
\end{itemize}

\subsection{Mathematical vs.\ Physical Grounding}

A crucial distinction that our audit surfaced: using real
mathematical equations does not make a mechanic physically
grounded.  Conformal geometry is real mathematics.  Applying a
conformal factor driven by an integration metric is not real
physics.  The realism table should distinguish:

\begin{itemize}
  \item \textbf{Mathematically grounded}: the equations are
    valid mathematics (correctly implemented).
  \item \textbf{Physically grounded}: the equations describe
    a relationship that holds in nature.
\end{itemize}

Our conformal curvature mechanic is mathematically grounded
but not physically grounded.  Failing to make this distinction
is, in our experience, the most common source of accidental
misrepresentation.

\subsection{Terminology Discipline}

We advocate for the following terminological practices:

\begin{itemize}
  \item Use ``integration metric'' or ``integration measure''
    rather than ``consciousness'' when describing the scalar
    quantity that drives coupling, unless the system explicitly
    claims to implement a specific theory of consciousness
    (e.g., IIT's $\Phi$) and the implementation can be validated
    against that theory's formal definitions.
  \item Avoid borrowing vocabulary from fundamental physics
    (``gravity,'' ``dimensions,'' ``curvature'') unless the
    mechanic faithfully implements the relevant physics.
    Metaphorical use of physics vocabulary should be flagged
    explicitly: ``we use the term `gravity' metaphorically
    to describe\ldots''
  \item When citing a theory as inspiration, specify what was
    omitted.  ``Inspired by the Free Energy Principle
    \citep{friston2010}, simplified to a scalar gradient
    (omitting the variational framework, generative model,
    and Markov blanket formalism)'' is more useful than
    ``based on the Free Energy Principle.''
\end{itemize}

\subsection{Emergent Behavior Disclaimers}

When reporting emergent behavior from \textsc{Speculative}
mechanics, papers should note that the behavior is an artifact
of the invented mechanic, not a prediction of any physical
theory.  This does not diminish the behavior's interest as a
simulation result.  It prevents the behavior from being
misinterpreted as evidence for a physical claim that was never
established.

\subsection{Versioned Realism Audits}

We recommend that realism audits be versioned and updated as
the system evolves.  Mechanics may move between tiers as new
science is published (a \textsc{Speculative} mechanic could
become \textsc{Inspired} if a relevant theory emerges) or as
implementations are revised (an \textsc{Inspired} mechanic
could become \textsc{Grounded} if the full source theory's
equations are implemented).


% ============================================================
\section{Related Work}
\label{sec:related}

The need for transparency in consciousness-related AI claims
has been raised from several directions.  \citet{chalmers1995}
established the conceptual distinction between ``easy'' and
``hard'' problems that makes integration-metric claims
philosophically fraught.  \citet{koch2019} has argued for
rigorous empirical criteria before attributing consciousness
to any system, biological or artificial.  \citet{bostrom2003}
raised the question of substrate-independence of consciousness
in simulation, but did not address the transparency of the
simulation's mechanics.

In the philosophy of simulation, \citet{winsberg2010} has
argued that models should be transparent about their
idealizations and approximations.  Our framework operationalizes
this principle for the specific case of integration-coupled
systems.

The IIT community has developed formal tools for computing
$\Phi$ \citep{tononi2004,oizumi2014,tononi2016}, and these
provide an example of \textsc{Grounded} practice: the theory
specifies its mathematics precisely enough that implementations
can be validated.  Our concern is with the many systems that
borrow IIT's vocabulary without implementing its formalism.

The FEP community has similarly developed formal tools
\citep{friston2010,parr2022}, but the gap between the full
FEP formalism and typical ``FEP-inspired'' implementations
in robotics and AI is substantial and often undisclosed.

To our knowledge, no prior work has proposed a systematic
classification framework for the realism of integration-physics
couplings, nor has any system published a full self-audit of
the kind we present here.


% ============================================================
\section{Limitations}
\label{sec:limitations}

Our framework has several limitations that should be
acknowledged.

First, the boundaries between tiers are judgment calls.  We
classified harmony fields as \textsc{Inspired} rather than
\textsc{Speculative} because the qualitative idea of fields
as integration substrates has a published advocate.  Others
might reasonably disagree.  The framework's value lies not in
producing unique classifications but in forcing the
classification conversation to happen.

Second, we have applied the framework only to our own engine.
We have not audited other systems, both because we lack access
to their internals and because the audit is most valuable as a
self-assessment tool.  We encourage other developers to
perform their own audits rather than relying on external
classification.

Third, the framework addresses only the \emph{physics} side of
integration-coupled systems.  It does not address whether the
integration metric itself is a valid measure of anything
philosophically interesting.  Whether $\Phi$ or any other
metric captures genuine consciousness is orthogonal to whether
the metric's coupling to physics is grounded, inspired, or
speculative.

Fourth, our audit was performed by the engine's developers.
Self-audits are subject to motivated reasoning.  We have tried
to err on the side of classifying mechanics as
\textsc{Speculative} when in doubt, but external audits would
be more rigorous.


% ============================================================
\section{Conclusion}
\label{sec:conclusion}

We built an agent simulation engine that couples an integration
metric to physical dynamics through nine distinct mechanics.
When we audited those mechanics for scientific grounding, we
found that only two rest on established physics, three draw
loose inspiration from real science, and four are our own
inventions dressed in physics vocabulary.

This finding did not diminish the engine.  The speculative
mechanics produce interesting emergent behavior.  But the audit
changed how we describe that behavior, how we document the
engine, and how we think about the relationship between
simulation and science.

The core lesson is simple: \textbf{transparency is a feature,
not a weakness}.  An engine that clearly labels its speculative
mechanics is more useful to researchers, safer for deployment,
and more honest to its users than one that presents invention
as established science.

We propose the \textsc{Grounded}/\textsc{Inspired}/\textsc{%
Speculative} framework as a minimal standard for any system
that couples integration metrics to simulated physics.  The
framework is deliberately simple---three tiers, clear criteria,
a single table.  Complexity is not the barrier to adoption;
cultural norms are.  If the field adopts a norm of disclosing
what is real and what is not, the quality of the science, the
safety of the systems, and the integrity of the philosophical
discourse will all benefit.

We have shown that performing this audit is tractable, that it
does not require removing speculative mechanics, and that it
improves both the engine's documentation and its developers'
understanding.  We encourage others to do the same.


% ============================================================
\section*{Acknowledgments}

The author thanks the open-source communities around IIT, the
Free Energy Principle, and Holochain for making the foundational
work accessible.  The Symtropy engine is released under the
AGPL-3.0-or-later license.


% ============================================================
\begin{thebibliography}{99}

\bibitem[B\'{e}rut et~al.(2012)]{berut2012}
B\'{e}rut, A., Arakelyan, A., Petrosyan, A., Ciliberto, S.,
  Dillenschneider, R., \& Lutz, E. (2012).
\newblock Experimental verification of {L}andauer's principle
  linking information and thermodynamics.
\newblock \emph{Nature}, 483(7388), 187--189.

\bibitem[Bostrom(2003)]{bostrom2003}
Bostrom, N. (2003).
\newblock Are we living in a computer simulation?
\newblock \emph{The Philosophical Quarterly}, 53(211), 243--255.

\bibitem[Carroll(2004)]{carroll2004}
Carroll, S.~M. (2004).
\newblock \emph{Spacetime and Geometry: An Introduction to
  General Relativity}.
\newblock Addison-Wesley.

\bibitem[Chalmers(1995)]{chalmers1995}
Chalmers, D.~J. (1995).
\newblock Facing up to the problem of consciousness.
\newblock \emph{Journal of Consciousness Studies}, 2(3),
  200--219.

\bibitem[Friston(2010)]{friston2010}
Friston, K. (2010).
\newblock The free-energy principle: A unified brain theory?
\newblock \emph{Nature Reviews Neuroscience}, 11(2), 127--138.

\bibitem[Kaufman et~al.(2014)]{kaufman2014}
Kaufman, M.~T., Churchland, M.~M., Ryu, S.~I., \&
  Shenoy, K.~V. (2014).
\newblock Cortical activity in the null space: Permitting
  preparation without movement.
\newblock \emph{Nature Neuroscience}, 17(3), 440--448.

\bibitem[Koch(2019)]{koch2019}
Koch, C. (2019).
\newblock \emph{The Feeling of Life Itself: Why Consciousness
  Is Widespread but Can't Be Computed}.
\newblock MIT Press.

\bibitem[Landauer(1961)]{landauer1961}
Landauer, R. (1961).
\newblock Irreversibility and heat generation in the computing
  process.
\newblock \emph{IBM Journal of Research and Development}, 5(3),
  183--191.

\bibitem[McFadden(2020)]{mcfadden2020}
McFadden, J. (2020).
\newblock Integrating information in the brain's {EM} field:
  The {CEMI} field theory of consciousness.
\newblock \emph{Neuroscience of Consciousness}, 2020(1),
  niaa016.

\bibitem[Oizumi et~al.(2014)]{oizumi2014}
Oizumi, M., Albantakis, L., \& Tononi, G. (2014).
\newblock From the phenomenology to the mechanisms of
  consciousness: {I}ntegrated {I}nformation {T}heory 3.0.
\newblock \emph{PLoS Computational Biology}, 10(5), e1003588.

\bibitem[Parr et~al.(2022)]{parr2022}
Parr, T., Pezzulo, G., \& Friston, K.~J. (2022).
\newblock \emph{Active Inference: The Free Energy Principle in
  Mind, Brain, and Behavior}.
\newblock MIT Press.

\bibitem[Randall \& Sundrum(1999)]{randall1999}
Randall, L. \& Sundrum, R. (1999).
\newblock Large mass hierarchy from a small extra dimension.
\newblock \emph{Physical Review Letters}, 83(17), 3370--3373.

\bibitem[Tononi(2004)]{tononi2004}
Tononi, G. (2004).
\newblock An information integration theory of consciousness.
\newblock \emph{BMC Neuroscience}, 5, 42.

\bibitem[Tononi et~al.(2016)]{tononi2016}
Tononi, G., Boly, M., Massimini, M., \& Koch, C. (2016).
\newblock Integrated information theory: From consciousness to
  its physical substrate.
\newblock \emph{Nature Reviews Neuroscience}, 17(7), 450--461.

\bibitem[Wald(1984)]{wald1984}
Wald, R.~M. (1984).
\newblock \emph{General Relativity}.
\newblock University of Chicago Press.

\bibitem[Winsberg(2010)]{winsberg2010}
Winsberg, E. (2010).
\newblock \emph{Science in the Age of Computer Simulation}.
\newblock University of Chicago Press.

\end{thebibliography}

\end{document}
